\chapter{Apertura del Tomo III}

\section{Función del tomo}
\begin{observation}
El Tomo III recibe piezas ya suficientemente preparadas en el canon intermedio y las trata en un nivel metodológico orientado a selección, demo, trazabilidad ampliada y posible traducción a claims.
\end{observation}

\begin{observation}
Su función no es refundar primitivas ni rehacer el canon operativo. Su tarea consiste en tomar material ya estabilizado y someterlo a criterios de admisión metodológica, uso auditado y proyección selectiva.
\end{observation}

\section{Relación con los tomos anteriores}
\begin{observation}
Tomo I fija el volumen fundacional principal del sistema: fundamento, blindaje, significado y núcleo estructural mínimo visible para sostener el resto del corpus.
\end{observation}

\begin{observation}
Los materiales técnicos del antiguo bloque Base Axiomática pueden preservarse como soporte interno de referencia, pero el fundamento público del sistema ya queda absorbido en Tomo I.
\end{observation}

\begin{observation}
El Tomo II fija el canon intermedio: transición, umbral, jerarquía, comparabilidad, configuración, familia, traducción, protocolización, trazabilidad y preparación para Tomo III.
\end{observation}

\begin{observation}
El Tomo III toma solo piezas que ya puedan entrar en tratamiento metodológico sin opacidad estructural inmediata.
\end{observation}

\section{Cláusula de no duplicación}
\begin{observation}
Nada de lo que entre al Tomo III debe refundar Tomo I ni duplicar el canon operativo intermedio de Tomo II salvo por referencia disciplinada estrictamente necesaria.
\end{observation}

\begin{observation}
Cuando una pieza dependa de bloques anteriores, esa dependencia debe mantenerse visible como herencia y no como reescritura encubierta del fundamento o del canon intermedio.
\end{observation}

\section{Régimen de lectura}
\begin{observation}
El Tomo III opera en un nivel metodológico: más selectivo que el canon intermedio y más exigente respecto de preparación, trazabilidad, protocolización y función demostrativa o curatorial.
\end{observation}

\begin{observation}
Su éxito editorial depende de no mezclarse ni con el laboratorio bruto ni con la refundación teórica. Debe funcionar como zona de traducción disciplinada entre piezas preparadas y usos auditables.
\end{observation}

\section{Cláusula de alcance}
\begin{observation}
Este tomo organiza admisión metodológica, demos seleccionadas, trazas ampliadas y posible traducción a claims. No reemplaza a Tomo I como volumen fundacional principal ni absorbe la función del canon operativo intermedio fijado en Tomo II.
\end{observation}

